\scanpage{6}

\diagram{\begin{tikzcd}[ampersand replacement=\&,row sep=small,column sep=small]
 \&s_2\arrow[dl,"u_2"']\&\\
 s_3\arrow[dr,"u_3"']\&\&s_1\arrow[ul,"u_1"']\\
 \&s_0\arrow[ur,"u_0"']\&
\end{tikzcd}}

\[
 \mu_0+\mu_2=\mu_1+\mu_3\qquad \mu_3=\mu_0-\mu_1+\mu_2
\]
\[\resizebox{\ifdim\width>\linewidth\linewidth\else\width\fi}{!}{$\displaystyle
\begin{aligned}
 &\deleted{\unsure{\Sigma(\lambda)_4=(\sigma_i)\in V^2|\Sigma(V)_4=(\sigma_i)\in V^2}_4}\\
 &\mathbb M_4\longrightarrow P^2
 &u_0\wedge u_2&=u_0\wedge(u_0+u_{\unsure{2|1}})\\
 &&&=\deleted{\unsure{+}}-u_0\wedge(u_1+u_3)\\
 &&&=-\underbrace{u_0\wedge u_1}_{\mu_1}+\underbrace{u_3\wedge u_0}_{\mu_0}
\end{aligned}
$}\]
\[\resizebox{\ifdim\width>\linewidth\linewidth\else\width\fi}{!}{$\displaystyle
(u_0,u_1,u_2,u_3)\longrightarrow
 \bigl(\underbrace{u_3\wedge u_0}_{\mu_0},
       \underbrace{u_0\wedge u_1}_{\mu_1},
       \underbrace{u_1\wedge u_2}_{\mu_2}\bigr).
$}\]
\[\resizebox{\ifdim\width>\linewidth\linewidth\else\width\fi}{!}{$\displaystyle
\begin{aligned}
 &\deleted{u_0\wedge u_2}\\
 u_0\wedge u_2&=\mu_0-\mu_1\\
 u_1\wedge u_3&=\mu_{\unsure{1|0}}-\mu_2\\
 u_2\wedge u_0&=\mu_2-\mu_3=\mu_1-\mu_0=-u_0\wedge u_2\\
 u_3\wedge u_1&=\mu_3-\mu_0=\mu_2-\mu_1=-u_1\wedge u_3
\end{aligned}
$}\]

$\mathbb M_4\longrightarrow P^2_{\mathbb Z}$ \deleted{injectif} mono (car la connaissance des $\mu_0,\mu_1,\mu_2$ implique celle des \deleted{$\mu_3$}, \unsure{donc|tous} les $u_i\wedge u_j$ \ldots)

épi ?

$\mu_0,\mu_1,\mu_2,\mu_3$ \deleted{\unsure{soit}} donnés, \qquad $\mu_0+\mu_2=\mu_1+\mu_3$, $\mu_i$ pas tous nuls

\unsure{OPS} $\mu_1\ne 0$

On prend $u_0,u_1$ comme base \inserted{\text{de }V}, $u_0\wedge u_1$ base \inserted{e} de $\bigwedge^2 V$
\[
 u_2=\alpha u_0+\beta u_1,\qquad u_3=-(u_0+u_1+u_2)
\]
\[\resizebox{\ifdim\width>\linewidth\linewidth\else\width\fi}{!}{$\displaystyle
\begin{aligned}
u_0\wedge u_1&=1.e & \mu'_1&=1\\
u_1\wedge u_2&=-\alpha e & \mu'_2&=-\alpha\\
u_2\wedge u_3&=(\deleted{\unsure{\alpha}}\beta-\deleted{\unsure{\beta}}\alpha)e
 & \underline{\mu'_3}&\underline{=\deleted{\unsure{\alpha}}\beta-\deleted{\unsure{\beta}}\alpha}\\
u_3\wedge u_0&=(1+\beta)e & \mu'_0&=1+\beta
\end{aligned}
\qquad
\begin{aligned}
u_2\wedge u_3&=-u_2\wedge u_0\unsure{-|+}u_2\wedge u_1\\
&=\beta\underbrace{u_0\wedge u_1}_{e}\deleted{\unsure{+}}-\alpha e\\
u_3\wedge u_0&=-u_1\wedge u_0-u_2\wedge u_0\\
&=\deleted{\unsure{1|V}}e+\beta e
\end{aligned}
$}\]
\[
 \alpha=-\frac{\mu_2}{\mu_1}\qquad
 \beta-\alpha=\frac{\mu_3}{\mu_1}\qquad \text{i.e.}\quad
 \beta=\frac{\mu_3}{\mu_1}+\alpha=\frac{\mu_3-\mu_2}{\mu_1}.
\]
\[\resizebox{\ifdim\width>\linewidth\linewidth\else\width\fi}{!}{$\displaystyle
\begin{aligned}
 u_2&=-\frac{\mu_2}{\mu_1}u_0+\frac{\mu_3-\mu_2}{\mu_1}u_1\\
 u_3&=\left(\frac{\mu_2}{\mu_1}-1\right)u_0+
 \underbrace{\frac{\mu_2-\mu_3-\mu_1}{\mu_1}}_{-\mu_0}u_1
\end{aligned}
\qquad
\begin{aligned}
 &\mathbb M_4\simeq P^2_{\mathbb Z}\\
 &u_2=0\Longleftrightarrow
 \left\{\begin{gathered}\deleted{\mu_1\ne0}\\ \boxed{\mu_2=0,\mu_3=0}\end{gathered}\right.\\
 &\hspace{2cm}\mu_0=\mu_1
\end{aligned}
$}\]

Les quadrilatères \unsure{épinglés}, \unsure{i.e.|ceux} avec un côté nul, correspondent aux quatre pts \inserted{a_i\ (i=0,1,\unsure{2|3},3)} de $P^2_{\mathbb Z}$ $(\deleted{\unsure{P^3}}\mu_0+\mu_2=\mu_1+\mu_3)\subset P^{\unsure{3|2}}$ $(\mu_0,\mu_1,\mu_2,\mu_3)$ correspondant à deux coordonnées \unsure{consec.|conséc.} $\mu_i$ nulles (les deux autres égales) \unsure{Il n'y a|Il n'y en a} pas $3$ dans \unsure{$a_i$|ces|ce} qui
% archive sheet number: 12
